\documentclass[11pt,a4paper,parskip=half-]{scrartcl}

% ---------- Encoding & Sprache ----------
\usepackage{polyglossia}
\setdefaultlanguage{german}

% ---------- Fonts (macOS-System) ----------
\usepackage{fontspec}
\setmainfont{Charter}[
  Numbers = OldStyle,
  Ligatures = TeX
]
\setsansfont{Avenir Next}[
  Scale = 0.92
]
\setmonofont{Menlo}[Scale = 0.78]

% ---------- Layout ----------
\usepackage[a4paper,margin=2.4cm,top=2.8cm,bottom=2.8cm]{geometry}
\usepackage{microtype}
\linespread{1.18}
\setlength{\parindent}{0pt}

% ---------- Tabellen ----------
\usepackage{array}
\usepackage{tabularx}
\usepackage{booktabs}
\usepackage{longtable}
\usepackage{makecell}
\renewcommand{\arraystretch}{1.25}

% Schmale Body-Spalten, lange Text-Spalten
\newcolumntype{L}[1]{>{\raggedright\arraybackslash}p{#1}}
\newcolumntype{C}[1]{>{\centering\arraybackslash}p{#1}}

% ---------- Farben ----------
\usepackage{xcolor}
\definecolor{accent}{HTML}{1F3A5F}     % gedämpftes Blau
\definecolor{accent2}{HTML}{8C2D2D}    % gedämpftes Rot
\definecolor{rulecolor}{HTML}{B0B5BB}
\definecolor{headbg}{HTML}{F2F4F7}
\definecolor{notebg}{HTML}{FBF7EE}
\definecolor{noteborder}{HTML}{D4B872}
\definecolor{newtag}{HTML}{2E6F40}
\definecolor{changedtag}{HTML}{B36B00}
\definecolor{tabletitlebg}{HTML}{1F3A5F}

% ---------- Kopf / Fuß ----------
\usepackage{scrlayer-scrpage}
\pagestyle{scrheadings}
\clearpairofpagestyles
\setkomafont{pagehead}{\sffamily\small\color{accent}}
\setkomafont{pagefoot}{\sffamily\footnotesize\color{accent}}
\ohead{Datenbankstruktur \textbar{} DGaO-Proceedings}
\ofoot*{\thepage}
\ifoot*{Stand: 26.\,Mai 2026}

% ---------- Section-Stil ----------
\usepackage[explicit]{titlesec}
\titleformat{\section}
  {\sffamily\Large\bfseries\color{accent}}
  {\thesection}{0.8em}{#1}
  [\vspace{-0.3em}{\color{rulecolor}\rule{\linewidth}{0.4pt}}]
\titleformat{\subsection}
  {\sffamily\large\bfseries\color{accent}}
  {}{0pt}{#1}
\titleformat{\subsubsection}
  {\sffamily\normalsize\bfseries\color{accent2}}
  {}{0pt}{#1}
\titlespacing*{\section}{0pt}{1.6em}{0.8em}
\titlespacing*{\subsection}{0pt}{1.4em}{0.5em}

% ---------- Hyperlinks (Inhaltsverz.) ----------
\usepackage[hidelinks,bookmarks=true,bookmarksopen=true,
            pdftitle={Datenbankstruktur DGaO-Proceedings},
            pdfauthor={Tobias Pruß}]{hyperref}

% ---------- Eigene Box für Hinweise ----------
\usepackage[breakable,skins]{tcolorbox}
\newtcolorbox{hinweis}{
  enhanced,
  colback=notebg,
  colframe=noteborder,
  boxrule=0.6pt,
  arc=2pt,
  left=10pt, right=10pt, top=6pt, bottom=6pt,
  breakable,
  before skip=8pt, after skip=10pt
}

% ---------- Tabellen-Kopf-Stil ----------
\newcommand{\tabhead}[1]{\textsf{\textbf{\small\color{accent}#1}}}

% ---------- Tag-Boxen (Neu/Geändert/Unverändert) ----------
\newcommand{\tagneu}{\colorbox{newtag!15}{\textcolor{newtag}{\sffamily\footnotesize\bfseries\,NEU\,}}}
\newcommand{\taggeandert}{\colorbox{changedtag!15}{\textcolor{changedtag}{\sffamily\footnotesize\bfseries\,GEÄNDERT\,}}}
\newcommand{\tagunv}{\textcolor{gray}{\sffamily\footnotesize\,unverändert\,}}

% ---------- Tabellen-Titelband ----------
\newcommand{\tabletitle}[2]{%
  \par\vspace{1.1em}%
  \noindent%
  \begin{tcolorbox}[
    enhanced,
    colback=tabletitlebg,
    colframe=tabletitlebg,
    arc=2pt,
    boxrule=0pt,
    left=10pt, right=10pt, top=4pt, bottom=4pt,
    before skip=0pt, after skip=6pt
  ]
    {\sffamily\large\bfseries\color{white}\texttt{#1}}\hfill{#2}
  \end{tcolorbox}%
  \par\nobreak%
}

% ============================================================
\begin{document}

% ---------- Titelseite ----------
\begin{titlepage}
  \centering
  \vspace*{2.5cm}
  {\sffamily\Huge\bfseries\color{accent} Datenbankstruktur\par}
  \vspace{0.4cm}
  {\sffamily\Large DGaO-Proceedings\par}
  \vspace{1.6cm}
  {\color{rulecolor}\rule{0.45\linewidth}{0.6pt}\par}
  \vspace{0.6cm}
  {\sffamily\large Schema nach der Konsolidierung\\[0.15cm]
   \normalsize von Autoren- und Institutionsdaten\par}
  \vspace{0.6cm}
  {\color{rulecolor}\rule{0.45\linewidth}{0.6pt}\par}
  \vfill
  \begin{tabular}{rl}
    \sffamily Erstellt von & Tobias Pruß \\
    \sffamily Stand        & 26.\,Mai 2026 \\
    \sffamily Status       & Entwurf zur Abstimmung \\
  \end{tabular}
  \vspace{1.5cm}
\end{titlepage}

% ---------- Inhalt ----------
\tableofcontents
\vspace{1.2em}

% ============================================================
\section{Überblick}

Dieses Dokument beschreibt die geplante Struktur der Datenbank,
die hinter der Webseite \emph{dgao-proceedings.de} steht. Es richtet sich an
Personen, die die Inhalte der Webseite pflegen oder weitergeben, und nicht
an Programmiererinnen oder Programmierer; technische Bezeichnungen werden
nur dort verwendet, wo sie wirklich nötig sind.

Die hier gezeigte Struktur ist der \textbf{geplante Endzustand}, nachdem die
seit 2004 angesammelten Schreibvarianten von Autorennamen und Instituten
zusammengeführt wurden. Sie weicht in mehreren Punkten von der heutigen
Datenbank ab, daher kennzeichnet jede Tabelle ihren Status:

\smallskip
\begin{tabular}{ll}
  \tagneu       & Tabelle existiert heute nicht und wird neu angelegt. \\
  \taggeandert  & Tabelle existiert, bekommt aber zusätzliche oder veränderte Felder. \\
  \tagunv       & Tabelle bleibt vollständig wie heute. \\
\end{tabular}

\subsection*{Wozu die Konsolidierung dient}

In den heutigen Daten gibt es dasselbe Institut in bis zu 16 verschiedenen
Schreibweisen, und dieselbe Person taucht oft unter mehreren Einträgen auf
(z.\,B.~\textit{C.~Pruß}, \textit{C.~Pruss}, \textit{Ch.~Pruß},
\textit{C.~Pruß*}). Das hat zwei Folgen:

\begin{itemize}\setlength{\itemsep}{2pt}
  \item Die \textbf{Suche auf der Webseite findet zu wenig}, weil sie
        Schreibvarianten nicht erkennt.
  \item Die Anzeige eines Autorenprofils zeigt nur die Beiträge, die
        unter exakt einer Schreibweise eingetragen wurden.
\end{itemize}

Die neue Struktur löst beide Probleme, ohne die Inhalte der bereits
veröffentlichten PDFs zu verändern: Die PDFs bleiben unangetastet,
die Datenbank pflegt jedoch zusätzlich eine saubere, einheitliche Sicht
auf Personen und Institute.

\subsection*{Wie die Tabellen zusammenhängen}

\begin{description}\setlength{\itemsep}{4pt}
  \item[Stammdaten der Tagungen] In \texttt{tagungen} ist jede DGaO-Tagung
        mit Nummer, Jahr und Ort hinterlegt. Die einzelnen Sitzungen
        (Vormittag-Block~A, Postersession, \dots) liegen in
        \texttt{sessions}.
  \item[Beiträge] Jeder Konferenzbeitrag (Vortrag oder Poster) ist eine
        Zeile in \texttt{papers}. Die Verknüpfung zu seinen Autorinnen und
        Autoren steht in \texttt{paper\_autoren}, die zu Stichwörtern in
        \texttt{paper\_keywords}.
  \item[Personen und Institute (neu)] Eine Person hat genau eine Zeile in
        \texttt{autoren} mit der offiziellen Schreibweise des Namens. Alle
        in den PDFs vorkommenden Schreibvarianten dieser Person stehen in
        \texttt{autor\_aliase}. Genauso hat ein Institut genau eine Zeile
        in \texttt{institutionen} und alle PDF-Schreibvarianten in
        \texttt{institut\_aliase}. Welche Person zu welchem Institut
        gehört (aktuell oder früher), regelt
        \texttt{autor\_institutionen}.
  \item[Weiteres] \texttt{news} steuert die Neuigkeiten auf der Startseite,
        \texttt{submissions} verwaltet die Einreichungen für neue Beiträge,
        \texttt{admin\_login\_attempts} dient der Sicherheit beim
        Administratorlogin.
\end{description}

\clearpage

% ============================================================
\section{Stammdaten der Tagungen}
% ============================================================

% ------------- tagungen -------------
\tabletitle{tagungen}{\tagunv}

\textit{Alle DGaO-Tagungen seit 2004, je eine Zeile pro Tagung.}

\begin{tabularx}{\linewidth}{L{3.6cm} L{1.9cm} L{3.0cm} X}
\toprule
\tabhead{Spalte} & \tabhead{Typ} & \tabhead{Pflicht} & \tabhead{Bedeutung} \\
\midrule
\texttt{nummer}              & Ganzzahl & ja, Primärschlüssel & Fortlaufende Tagungsnummer (z.\,B.~125 für die DGaO~2024). \\
\texttt{jahr}                & Ganzzahl & ja & Jahr der Tagung. \\
\texttt{ort}                 & Text     & nein & Veranstaltungsort. \\
\texttt{datum\_von}          & Text     & nein & Erster Tagungstag. \\
\texttt{datum\_bis}          & Text     & nein & Letzter Tagungstag. \\
\texttt{vorlage\_phase\_aktiv} & Ganzzahl (0/1) & ja & Steuert, ob die Einreichungsphase aktiv ist. \\
\texttt{einreichungsfrist}   & Text     & nein & Letztes Datum für Einreichungen. \\
\bottomrule
\end{tabularx}

% ------------- sessions -------------
\tabletitle{sessions}{\tagunv}

\textit{Die einzelnen Sitzungsblöcke einer Tagung (z.\,B.~Block~A, Postersession).
Eine Sitzung ohne zugeordneten Beitrag wird nicht gespeichert.}

\begin{tabularx}{\linewidth}{L{3.6cm} L{1.9cm} L{3.0cm} X}
\toprule
\tabhead{Spalte} & \tabhead{Typ} & \tabhead{Pflicht} & \tabhead{Bedeutung} \\
\midrule
\texttt{id}              & Ganzzahl & ja, Primärschlüssel & Eindeutige interne Kennung der Sitzung. \\
\texttt{tagung\_nummer}  & Ganzzahl & ja, verweist auf \texttt{tagungen} & Zu welcher Tagung gehört die Sitzung. \\
\texttt{titel}           & Text     & ja & Titel der Sitzung (z.\,B.~\textit{Holografie I}). \\
\texttt{saal}            & Text     & nein & Saalbezeichnung (z.\,B.~\textit{A}). \\
\texttt{sortorder}       & Ganzzahl & ja & Reihenfolge im Programm. \\
\texttt{datum}           & Text     & nein & Datum der Sitzung. \\
\texttt{zeit\_von}       & Text     & nein & Beginn (Uhrzeit). \\
\texttt{zeit\_bis}       & Text     & nein & Ende (Uhrzeit). \\
\bottomrule
\end{tabularx}

% ============================================================
\section{Beiträge}
% ============================================================

% ------------- papers -------------
\tabletitle{papers}{\tagunv}

\textit{Ein Konferenzbeitrag pro Zeile (Vortrag oder Poster).
Die hier gespeicherten Texte stammen direkt aus den PDFs und werden nicht
verändert.}

\begin{tabularx}{\linewidth}{L{3.6cm} L{1.9cm} L{3.0cm} X}
\toprule
\tabhead{Spalte} & \tabhead{Typ} & \tabhead{Pflicht} & \tabhead{Bedeutung} \\
\midrule
\texttt{id}               & Text     & ja, Primärschlüssel & Eindeutige Kennung des Beitrags. \\
\texttt{tagung\_nummer}   & Ganzzahl & ja, verweist auf \texttt{tagungen} & Zu welcher Tagung gehört der Beitrag. \\
\texttt{code}             & Text     & ja & Programmcode (z.\,B.~\textit{A1}, \textit{P12}). \\
\texttt{typ}              & Text     & ja & Beitragstyp (z.\,B.~\textit{Vortrag}, \textit{Poster}). \\
\texttt{titel}            & Text     & ja & Titel des Beitrags. \\
\texttt{autoren\_text}    & Text     & ja & Autorenliste, exakt wie im PDF gedruckt. \\
\texttt{hauptautor}       & Text     & nein & Name des Hauptautors als Anzeigetext. \\
\texttt{abstract\_text}   & Text     & nein & Kurzfassung. \\
\texttt{zeit}             & Text     & nein & Vortragszeit (falls hinterlegt). \\
\texttt{raum}             & Text     & nein & Raum im Tagungsprogramm. \\
\texttt{datum}            & Text     & nein & Datum des Beitrags. \\
\texttt{affiliationen}    & Text     & nein & Institutsangaben, exakt wie im PDF gedruckt. \\
\texttt{kontakt\_email}   & Text     & nein & Kontakt-E-Mail-Adresse (sofern angegeben). \\
\texttt{pdf\_dateiname}   & Text     & nein & Dateiname des PDFs auf dem Server. \\
\texttt{hat\_pdf}         & Ganzzahl (0/1) & ja & Kennzeichen, ob ein PDF hinterlegt ist. \\
\texttt{alte\_abstract\_id} & Ganzzahl & nein & Bezug zur alten Datenstruktur (nur historisch). \\
\texttt{session\_id}      & Ganzzahl & nein, verweist auf \texttt{sessions} & Zugehörige Sitzung im Programm. \\
\bottomrule
\end{tabularx}

\begin{hinweis}
Die Felder \texttt{autoren\_text} und \texttt{affiliationen} enthalten die
Rohtexte aus den PDFs und bleiben unverändert. Die strukturierten,
bereinigten Informationen pflegt die Datenbank parallel über die Tabellen
\texttt{autoren}, \texttt{autor\_aliase}, \texttt{institutionen} und
\texttt{institut\_aliase}.
\end{hinweis}

% ------------- paper_autoren -------------
\tabletitle{paper\_autoren}{\tagunv}

\textit{Verknüpfung zwischen Beiträgen und Autoren. Eine Zeile pro
Autorenrolle eines Beitrags.}

\begin{tabularx}{\linewidth}{L{3.6cm} L{1.9cm} L{3.0cm} X}
\toprule
\tabhead{Spalte} & \tabhead{Typ} & \tabhead{Pflicht} & \tabhead{Bedeutung} \\
\midrule
\texttt{paper\_id}       & Text     & ja, Primärschlüssel (gemeinsam) & Verweis auf den Beitrag in \texttt{papers}. \\
\texttt{autor\_id}       & Ganzzahl & ja, Primärschlüssel (gemeinsam) & Verweis auf die Person in \texttt{autoren}. \\
\texttt{position}        & Ganzzahl & ja & Reihenfolge der Autorin oder des Autors. \\
\texttt{ist\_hauptautor} & Ganzzahl (0/1) & nein & Kennzeichen für Hauptautorenrolle. \\
\bottomrule
\end{tabularx}

\clearpage

% ============================================================
\section{Personen und Institute}
% ============================================================

% ------------- autoren -------------
\tabletitle{autoren}{\taggeandert}

\textit{Eine Zeile pro Person mit der offiziellen Schreibweise des Namens.
Alle weiteren Schreibvarianten aus den PDFs werden separat in
\texttt{autor\_aliase} geführt.}

\begin{tabularx}{\linewidth}{L{3.6cm} L{1.9cm} L{3.0cm} X}
\toprule
\tabhead{Spalte} & \tabhead{Typ} & \tabhead{Pflicht} & \tabhead{Bedeutung} \\
\midrule
\texttt{id}        & Ganzzahl & ja, Primärschlüssel & Eindeutige interne Kennung der Person. \\
\texttt{vorname}   & Text     & ja, gemeinsam mit \texttt{nachname} eindeutig & Vorname oder Initialen in offizieller Schreibweise. \\
\texttt{nachname}  & Text     & ja, gemeinsam mit \texttt{vorname} eindeutig & Nachname in offizieller Schreibweise. \\
\texttt{orcid\_id} & Text     & nein & ORCID-Kennung der Person (optional, derzeit nicht gepflegt). \\
\bottomrule
\end{tabularx}

\begin{hinweis}
\textbf{Was bedeutet „offizielle Schreibweise"?} Eine Person hat
genau eine Zeile in dieser Tabelle, mit der Schreibweise, die für die
Anzeige verwendet wird (z.\,B.~\textit{C.~Pruß}). Fußnoten-Sterne aus den
PDFs sowie alternative Schreibweisen werden nicht hier, sondern in
\texttt{autor\_aliase} gespeichert.
\end{hinweis}

% ------------- autor_aliase -------------
\tabletitle{autor\_aliase}{\tagneu}

\textit{Alle Schreibvarianten eines Namens, mit denen eine Person je in
einem PDF aufgetreten ist. Erlaubt es der Suche, eine Person unabhängig
von der konkreten Schreibweise zu finden.}

\begin{tabularx}{\linewidth}{L{3.6cm} L{1.9cm} L{3.0cm} X}
\toprule
\tabhead{Spalte} & \tabhead{Typ} & \tabhead{Pflicht} & \tabhead{Bedeutung} \\
\midrule
\texttt{id}          & Ganzzahl & ja, Primärschlüssel & Eindeutige Kennung des Alias-Eintrags. \\
\texttt{autor\_id}   & Ganzzahl & ja, verweist auf \texttt{autoren} & Zu welcher Person gehört die Schreibvariante. \\
\texttt{alias\_text} & Text     & ja & Schreibvariante so, wie sie im PDF stand (ohne Fußnoten-Sterne). \\
\texttt{alias\_norm} & Text     & ja, gemeinsam mit \texttt{autor\_id} eindeutig & Vereinheitlichte Form für die Suche (Kleinschreibung, ohne Sonderzeichen). \\
\texttt{created\_at} & Text     & ja & Zeitpunkt, zu dem die Variante angelegt wurde. \\
\bottomrule
\end{tabularx}

\begin{hinweis}
\textbf{Beispiel.} Die Person \textit{C.~Pruß} hat in Wirklichkeit unter
fünf verschiedenen Schreibweisen publiziert. Nach der Konsolidierung gibt
es eine Zeile in \texttt{autoren} und mehrere Zeilen in
\texttt{autor\_aliase}, eine pro Schreibvariante. Die Suche findet die
Person unabhängig davon, welche Schreibweise eingegeben wird.
\end{hinweis}

% ------------- institutionen -------------
\tabletitle{institutionen}{\tagneu}

\textit{Eine Zeile pro Institut mit der offiziellen deutschen und
englischen Bezeichnung, dem Kürzel und der zugehörigen Hochschule oder
Forschungseinrichtung.}

\begin{tabularx}{\linewidth}{L{3.6cm} L{1.9cm} L{3.0cm} X}
\toprule
\tabhead{Spalte} & \tabhead{Typ} & \tabhead{Pflicht} & \tabhead{Bedeutung} \\
\midrule
\texttt{id}            & Ganzzahl & ja, Primärschlüssel & Eindeutige interne Kennung des Instituts. \\
\texttt{name\_de}      & Text     & ja & Offizielle deutsche Bezeichnung. \\
\texttt{name\_en}      & Text     & nein & Offizielle englische Bezeichnung. Bei leerem Wert wird die deutsche Bezeichnung angezeigt. \\
\texttt{kuerzel}       & Text     & nein & Übliche Abkürzung (z.\,B.~\textit{ITO}, \textit{PTB}). \\
\texttt{universitaet}  & Text     & nein & Übergeordnete Hochschule oder Einrichtung (z.\,B.~\textit{Universität Stuttgart}). \\
\texttt{ort}           & Text     & nein & Sitz des Instituts. \\
\texttt{land}          & Text     & ja, Vorgabe \textit{DE} & ISO-Länderkürzel. \\
\texttt{ror\_id}       & Text     & nein & Kennung im internationalen ROR-Register (optional). \\
\texttt{created\_at}   & Text     & ja & Zeitpunkt der Anlage. \\
\bottomrule
\end{tabularx}

% ------------- institut_aliase -------------
\tabletitle{institut\_aliase}{\tagneu}

\textit{Alle Schreibvarianten eines Instituts, mit denen es je in einem PDF
aufgetreten ist.}

\begin{tabularx}{\linewidth}{L{3.6cm} L{1.9cm} L{3.0cm} X}
\toprule
\tabhead{Spalte} & \tabhead{Typ} & \tabhead{Pflicht} & \tabhead{Bedeutung} \\
\midrule
\texttt{id}             & Ganzzahl & ja, Primärschlüssel & Eindeutige Kennung des Alias-Eintrags. \\
\texttt{institut\_id}   & Ganzzahl & ja, verweist auf \texttt{institutionen} & Zu welchem Institut gehört die Schreibvariante. \\
\texttt{alias\_text}    & Text     & ja & Schreibvariante so, wie sie im PDF stand. \\
\texttt{alias\_norm}    & Text     & ja, gemeinsam mit \texttt{institut\_id} eindeutig & Vereinheitlichte Form für die Suche. \\
\texttt{created\_at}    & Text     & ja & Zeitpunkt der Anlage. \\
\bottomrule
\end{tabularx}

% ------------- autor_institutionen -------------
\tabletitle{autor\_institutionen}{\tagneu}

\textit{Verknüpfung zwischen Personen und Instituten. Eine Zeile pro
Personen-Institut-Beziehung. Eine Person kann zu mehreren Instituten
gehört haben; genau eines davon ist als „aktuell" markiert.}

\begin{tabularx}{\linewidth}{L{3.6cm} L{1.9cm} L{3.0cm} X}
\toprule
\tabhead{Spalte} & \tabhead{Typ} & \tabhead{Pflicht} & \tabhead{Bedeutung} \\
\midrule
\texttt{autor\_id}    & Ganzzahl & ja, Primärschlüssel (gemeinsam) & Verweis auf die Person in \texttt{autoren}. \\
\texttt{institut\_id} & Ganzzahl & ja, Primärschlüssel (gemeinsam) & Verweis auf das Institut in \texttt{institutionen}. \\
\texttt{ist\_aktuell} & Ganzzahl (0/1) & ja, Vorgabe~\texttt{0} & Kennzeichen, ob es sich um das aktuelle Institut der Person handelt. \\
\bottomrule
\end{tabularx}

\begin{hinweis}
\textbf{Wie wird das aktuelle Institut bestimmt?} Die Datenbank betrachtet
das Institut, das in der zeitlich jüngsten Tagung mit einem Beitrag dieser
Person verknüpft war, automatisch als das aktuelle. Alle weiteren
Institute werden als frühere Zugehörigkeiten geführt.
\end{hinweis}

% ------------- autor_id_redirects -------------
\tabletitle{autor\_id\_redirects}{\tagneu}

\textit{Hilfstabelle: Wenn mehrere Datensätze zur selben Person
zusammengeführt werden, merkt sich die Datenbank hier, dass die alten
Kennungen auf die neue Kennung verweisen. Damit bleiben bereits extern
verlinkte Autorenseiten erreichbar.}

\begin{tabularx}{\linewidth}{L{3.6cm} L{1.9cm} L{3.0cm} X}
\toprule
\tabhead{Spalte} & \tabhead{Typ} & \tabhead{Pflicht} & \tabhead{Bedeutung} \\
\midrule
\texttt{alte\_id}    & Ganzzahl & ja, Primärschlüssel & Frühere Kennung der Person (vor der Zusammenführung). \\
\texttt{neue\_id}    & Ganzzahl & ja, verweist auf \texttt{autoren} & Aktuelle, gültige Kennung. \\
\texttt{merged\_at}  & Text     & ja & Zeitpunkt der Zusammenführung. \\
\bottomrule
\end{tabularx}

\clearpage

% ============================================================
\section{Stichwörter}
% ============================================================

% ------------- keywords -------------
\tabletitle{keywords}{\tagunv}

\textit{Liste der Stichwörter, mit denen Beiträge versehen sind.}

\begin{tabularx}{\linewidth}{L{3.6cm} L{1.9cm} L{3.0cm} X}
\toprule
\tabhead{Spalte} & \tabhead{Typ} & \tabhead{Pflicht} & \tabhead{Bedeutung} \\
\midrule
\texttt{id}      & Ganzzahl & ja, Primärschlüssel & Eindeutige interne Kennung. \\
\texttt{keyword} & Text     & ja, eindeutig & Das Stichwort selbst. \\
\bottomrule
\end{tabularx}

% ------------- paper_keywords -------------
\tabletitle{paper\_keywords}{\tagunv}

\textit{Verknüpfung zwischen Beiträgen und Stichwörtern.}

\begin{tabularx}{\linewidth}{L{3.6cm} L{1.9cm} L{3.0cm} X}
\toprule
\tabhead{Spalte} & \tabhead{Typ} & \tabhead{Pflicht} & \tabhead{Bedeutung} \\
\midrule
\texttt{paper\_id}    & Text     & ja, Primärschlüssel (gemeinsam) & Verweis auf den Beitrag. \\
\texttt{keyword\_id}  & Ganzzahl & ja, Primärschlüssel (gemeinsam) & Verweis auf das Stichwort. \\
\bottomrule
\end{tabularx}

% ============================================================
\section{Redaktionelle Inhalte}
% ============================================================

% ------------- news -------------
\tabletitle{news}{\tagunv}

\textit{Neuigkeiten auf der Startseite. Manche werden automatisch erzeugt
(z.\,B.~bei Beginn einer Einreichungsphase), andere pflegt die
Administration manuell.}

\begin{tabularx}{\linewidth}{L{3.6cm} L{1.9cm} L{3.0cm} X}
\toprule
\tabhead{Spalte} & \tabhead{Typ} & \tabhead{Pflicht} & \tabhead{Bedeutung} \\
\midrule
\texttt{id}              & Ganzzahl & ja, Primärschlüssel & Eindeutige Kennung. \\
\texttt{source}          & Text     & ja & \textit{auto} oder \textit{manual}. \\
\texttt{trigger\_key}    & Text     & nein & Auslöser für automatische Neuigkeiten. \\
\texttt{tagung\_nummer}  & Ganzzahl & nein, verweist auf \texttt{tagungen} & Zugehörige Tagung (falls vorhanden). \\
\texttt{display\_date}   & Text     & ja & Auf der Startseite angezeigtes Datum. \\
\texttt{title\_de}       & Text     & ja & Deutscher Titel. \\
\texttt{title\_en}       & Text     & ja & Englischer Titel. \\
\texttt{body\_de}        & Text     & ja & Deutscher Fließtext. \\
\texttt{body\_en}        & Text     & ja & Englischer Fließtext. \\
\texttt{link\_url}       & Text     & nein & Zielseite des Eintrags. \\
\texttt{is\_active}      & Ganzzahl (0/1) & ja & Kennzeichen, ob die Neuigkeit sichtbar ist. \\
\texttt{sort\_weight}    & Ganzzahl & ja & Sortiergewicht (höhere Werte = oben gepinnt). \\
\texttt{created\_at}     & Text     & ja & Zeitpunkt der Anlage. \\
\texttt{updated\_at}     & Text     & ja & Zeitpunkt der letzten Änderung. \\
\bottomrule
\end{tabularx}

% ============================================================
\section{Einreichungen und Sicherheit}
% ============================================================

% ------------- submissions -------------
\tabletitle{submissions}{\tagunv}

\textit{Verwaltet die Einreichungen von Manuskripten für eine
Tagung. Jede Einreichung hat ein eindeutiges Token (Link für die
Einreichende oder den Einreichenden) und einen Status.}

\begin{tabularx}{\linewidth}{L{3.6cm} L{1.9cm} L{3.0cm} X}
\toprule
\tabhead{Spalte} & \tabhead{Typ} & \tabhead{Pflicht} & \tabhead{Bedeutung} \\
\midrule
\texttt{token}              & Text     & ja, Primärschlüssel & Geheimer Zugangsschlüssel für die Einreichende oder den Einreichenden. \\
\texttt{paper\_id}          & Text     & ja, verweist auf \texttt{papers} & Zugehöriger Beitrag. \\
\texttt{status}             & Text     & ja, Vorgabe~\textit{pending} & \textit{pending}, \textit{approved}, \textit{rejected} oder \textit{expired}. \\
\texttt{uploader\_email}    & Text     & ja & E-Mail-Adresse der einreichenden Person. \\
\texttt{filename\_original} & Text     & nein & Ursprünglicher Dateiname. \\
\texttt{filename\_stored}   & Text     & nein & Dateiname im Speicher. \\
\texttt{file\_size}         & Ganzzahl & nein & Dateigröße in Bytes. \\
\texttt{requested\_at}      & Text     & ja & Zeitpunkt der Anforderung. \\
\texttt{uploaded\_at}       & Text     & nein & Zeitpunkt des Hochladens. \\
\texttt{decided\_at}        & Text     & nein & Zeitpunkt der Entscheidung. \\
\texttt{decided\_by}        & Text     & nein & Wer die Entscheidung getroffen hat. \\
\texttt{reviewer\_note}     & Text     & nein & Notiz der prüfenden Person. \\
\texttt{expires\_at}        & Text     & ja & Ablaufdatum des Tokens. \\
\bottomrule
\end{tabularx}

% ------------- admin_login_attempts -------------
\tabletitle{admin\_login\_attempts}{\tagunv}

\textit{Protokoll der Anmeldeversuche im Administrationsbereich,
ausschließlich zur Begrenzung von missbräuchlichen Login-Versuchen.}

\begin{tabularx}{\linewidth}{L{3.6cm} L{1.9cm} L{3.0cm} X}
\toprule
\tabhead{Spalte} & \tabhead{Typ} & \tabhead{Pflicht} & \tabhead{Bedeutung} \\
\midrule
\texttt{id}  & Ganzzahl & ja, Primärschlüssel & Laufende Nummer. \\
\texttt{ip}  & Text     & ja & IP-Adresse, von der der Versuch ausging. \\
\texttt{ts}  & Ganzzahl & ja & Zeitstempel des Versuchs. \\
\bottomrule
\end{tabularx}

\clearpage

% ============================================================
\section{Zusammenfassung der Änderungen}
% ============================================================

Die folgende Übersicht zeigt, was sich gegenüber der heutigen Datenbank
ändert.

\begin{tabularx}{\linewidth}{L{4.6cm} L{2.6cm} X}
\toprule
\tabhead{Tabelle} & \tabhead{Status} & \tabhead{Was ändert sich} \\
\midrule
\texttt{tagungen}                & \tagunv      & keine Änderung \\
\texttt{sessions}                & \tagunv      & keine Änderung \\
\texttt{papers}                  & \tagunv      & keine Änderung \\
\texttt{paper\_autoren}          & \tagunv      & keine Änderung \\
\texttt{autoren}                 & \taggeandert & Spalte für ORCID-Kennung ergänzt; alte Spalte \texttt{affiliation} entfällt (wird durch \texttt{autor\_institutionen} ersetzt). \\
\texttt{autor\_aliase}           & \tagneu      & speichert alle Schreibvarianten eines Namens \\
\texttt{institutionen}           & \tagneu      & einheitliche Institutsdaten in Deutsch und Englisch \\
\texttt{institut\_aliase}        & \tagneu      & speichert alle Schreibvarianten eines Instituts \\
\texttt{autor\_institutionen}    & \tagneu      & verknüpft Personen mit aktuellen und früheren Instituten \\
\texttt{autor\_id\_redirects}    & \tagneu      & bewahrt extern verlinkte alte Autoren-Adressen \\
\texttt{keywords}                & \tagunv      & keine Änderung \\
\texttt{paper\_keywords}         & \tagunv      & keine Änderung \\
\texttt{news}                    & \tagunv      & keine Änderung \\
\texttt{submissions}             & \tagunv      & keine Änderung \\
\texttt{admin\_login\_attempts}  & \tagunv      & keine Änderung \\
\bottomrule
\end{tabularx}

\vspace{1.2em}

\begin{hinweis}
\textbf{Was bleibt unverändert?} Die hochgeladenen PDFs werden in keinem
Schritt verändert. Auch die Felder \texttt{autoren\_text} und
\texttt{affiliationen} in \texttt{papers}, die die Originaltexte aus den
PDFs enthalten, bleiben so wie sie sind. Die Konsolidierung findet
ausschließlich in der dahinter liegenden Datenbank statt.
\end{hinweis}

\end{document}
